%\input{../preamble}

%\usepackage[refpage]{nomencl}
%\makenomenclature

%\begin{document}

%\printnomenclature

\subsection{Groupoids as directed graphs}

\begin{definition}[\cite{MR1724106}]
A \emph{directed graph} is a tuple $\G=(\G[0],\G[1],\so,\ra)$, where $\mathcal{G}^{(0)}$ is a set whose elements are called \emph{vertices}, $\mathcal{G}^{(1)}$ is a set whose elements are called \emph{edges} or \emph{arrows}, and $\mathfrak{s}$ and $\mathfrak{r}$ are maps from $\mathcal{G}^{(1)}$ to $\mathcal{G}^{(0)}$, called respectively the \emph{source} and \emph{range} maps.
\end{definition}

Directed graphs can be depicted as points, representing the vertices, and arrows, representing the edges, pointing from its \emph{source} to its \emph{range}. 
\[\begin{tikzcd}
\mathfrak{s}(a)\arrow[rr,"a"]& &\mathfrak{r}(a)
\end{tikzcd}\qquad
\begin{tikzcd}[row sep=small]
\bullet\arrow[dr,shift left]\arrow[dr,shift right]\arrow[rr, bend right=15]&&\arrow[ll, bend right=15]\bullet\arrow[loop right]\\
&\bullet&\end{tikzcd}\]
It should be noted that the directed graphs we consider here are sometimes called \emph{multigraphs} or \emph{pseudographs}, since we allow loops and multiple edges between two vertices.

\begin{example}\label{examplegraphtriangle}
Let $X$ be a set and $R$ any binary relation on $X$, i.e., a subset of $X^2$. Define a directed graph $\G$ by letting $\G[0]=X$, $\G[1]=R$, $\mathfrak{s}(y,x)=x$ and $\mathfrak{r}(y,x)=y$ for all $(y,x)\in R$.

As a sub-example, let $X=\left\{a,b,c\right\}$ and $R=\left\{(b,a),(c,b),(a,c)\right\}$. The associated graph can be depicted as
\[\begin{tikzcd}
&b\arrow[rd]&\\
a\arrow[ur]&&c\arrow[ll]\end{tikzcd}\]
\end{example}

\begin{example}\label{examplegraphintegers}
Let $\mathcal{G}^{(0)}=\mathcal{G}^{(1)}=\mathbb{Z}$, $\mathfrak{s}(n)=n$ and $\mathfrak{r}(n)=n+1$. This directed graph can be depicted as
\[\begin{tikzcd}[column sep=small]
\cdots\arrow[r]&-2\arrow[r]&-1\arrow[r]&0\arrow[r]&1\arrow[r]&1\arrow[r]&\cdots\arrow[r]&n\arrow[r]&n+1\arrow[r]&\cdots
\end{tikzcd}\]
\end{example}

\begin{example}\label{examplegraphdynamicalsystem}
Generally, a discrete-time dynamical system consists of a set $X$ and a collection $\mathcal{F}$ of partial functions of $X$, that is, functions $f:\dom(f)\to X$ where $\dom(f)\subseteq X$. Given such data, we define a directed graph $\G$ by letting $\G[0]=X$, $\G[1]=\left\{(f,x):f\in\mathcal{F},x\in\dom(f)\right\}$, $\so(f,x)=x$ and $\ra(f,x)=f(x)$.

This directed graph $\G$ is depicted by drawing elements of $X$ as points and arrows from points of $X$ to their image under elements of $\mathcal{F}$, whenever defined.

As a sub-example, let $X=\mathbb{Z}$, and $f_i(n)=n+i$ for $i=0,1,2$. The associated graph is depicted as
\[\begin{tikzcd}
\cdots\arrow[r]\arrow[rr,bend right]
	&-2\arrow[r]\arrow[rr,bend right,crossing over]\arrow[loop above]
		&-1\arrow[r]\arrow[rr,bend right,crossing over]\arrow[loop above]
			&0\arrow[r]\arrow[rr,bend right,crossing over]\arrow[loop above]
				&1\arrow[r]\arrow[rr,bend right,crossing over]\arrow[loop above]
					& 2\arrow[r]\arrow[loop above]&\cdots\end{tikzcd}\]
\end{example}

In order to obtain a suitable notion for the composition of arrows, we need to introduce the concept of a \emph{category}.

\begin{definition}\nomenclature[G]{$\G_x$, $\G^y$, $\G_x^y$}{Arrows with specified source and range}
Let $\G=(\G[0],\G[1],\so,\ra)$ be a directed graph. Given vertices $x,y\in\G[0]$, denote by
\begin{itemize}
\item $\G_x=\mathfrak{s}^{-1}(x)$, the set of arrows leaving from $x$;
\item $\G^y=\mathfrak{r}^{-1}(y)$, the set of arrows arriving at $y$;
\item $\G_x^y=\G^y\cap\G_x$, the set of arrows from $x$ to $y$.
\end{itemize}
\end{definition}

\begin{definition}%[\cite{MR1712872}]
A \emph{small category} consists of a directed graph $\mathcal{C}=(\mathcal{C}^{(0)},\mathcal{C}^{(1)},\mathfrak{s},\mathfrak{r})$ together with the set $\mathcal{C}^{(2)}=\left\{(a,b)\in\mathcal{C}^{(1)}\times\mathcal{C}^{(1)}:\mathfrak{s}(a)=\mathfrak{r}(b)\right\}$ and a \emph{composition} map $\mathcal{C}^{(2)}\to \mathcal{C}^{(1)}$, $(a,b)\mapsto ab$, satisfying the following properties:
\begin{enumerate}
%\item[(i)] If $(a,b)\in\mathcal{C}^{(2)}$, then $\mathfrak{s}(ab)=\mathfrak{s}(b)$ and $\mathfrak{r}(ab)=\mathfrak{r}(a)$;
%\item[(ii)] (Associativity) If $(a,b)$ and $(b,c)\in\mathcal{C}^{(2)}$, then $a(bc)=(ab)c$. (Note that both sides of this equality make sense by (i).)
%\item[(iii)] (Identity) For every $x\in\mathcal{C}^{(0)}$, there exists $\operatorname{id}_x\in \mathcal{C}_x^x$ such that if $a\in\mathcal{C}_x$ and $b\in\mathcal{C}^x$, then $a\operatorname{id}_x=a$ and $\operatorname{id}_xb=b$.
\item[(i)] If $(a,b)$ and $(b,c)\in\mathcal{C}^{(2)}$, then $(a,bc),(ab,c)\in\mathcal{C}^{(2)}$ and $a(bc)=(ab)c$.
\item[(ii)] For every $x\in\mathcal{C}^{(0)}$, there exists $\operatorname{id}_x\in \mathcal{C}_x^x$ such that if $a\in\mathcal{C}_x$ and $b\in\mathcal{C}^ x$, then $a\operatorname{id}_x=a$ and $\operatorname{id}_xb=b$.
\end{enumerate}
The set $\mathcal{C}^{(2)}$ is called the set of \emph{composable pairs}.\index{Composable}
\end{definition}

General categories are defined in the same manner, however we allow $\mathcal{C}^{(0)}$ and $\mathcal{C}^{(1)}$ to be general classes instead of sets, however we will restrict our study to sets.

Note that $\mathcal{C}^{(2)}=\bigcup_{x\in\mathcal{C}^{(0)}}\mathcal{C}_x\times\mathcal{C}^x$, and so the pair $(a,b)$ of arrows in a category $\mathcal{C}$ is \emph{composable} if they meet ``tip-to-tail'':
\[\begin{tikzcd}[row sep=tiny]
\underset{\so(b)}{\bullet}\arrow[r,"b"]\arrow[rr,bend right,"ab",swap]&\underset{\ra(b)=\so(a)}{\bullet}\arrow[r,"a"]&\underset{\ra(a)}{\bullet}
\end{tikzcd}\]

The arrow $\operatorname{id}_x\in\mathcal{C}_x^x$ in item (ii) is unique for each $x\in\mathcal{C}^{(0)}$, and it is called the \emph{unit} at $x$. It is customary to identify each $x\in\mathcal{C}^{(0)}$ with its respective unit, thus viewing $\mathcal{C}^{(0)}$ as a subset of $\mathcal{C}^{(1)}$. This allows us, moreover, to identify a category with its set of arrows.

Moreover, from item (i) we obtain, whenever $(a,b)\in\mathcal{C}^{(2)}$, $(ab,\id_{\so(b)})\in\mathcal{C}^{(2)}$, and so $\so(ab)=\ra(\id_{\so(b)})=\so(b)$ and similarly $\ra(ab)=\ra(a)$.

\begin{definition}\index{Morphism (of groupoids)}\index{Isomorphism (of groupoids)}
A \emph{functor} between categories $\mathcal{C}$ and $\mathcal{D}$ is a map $\phi:\mathcal{C}\to\mathcal{D}$ satisfying:
\begin{enumerate}
\item[(i)] $\phi(\mathcal{C}^{(0)})\subseteq\mathcal{D}^{(0)}$.
\item[(ii)] If $(a,b)\in\mathcal{C}^{(2)}$, then $(\phi(a),\phi(b))\in\mathcal{D}^{(2)}$ and $\phi(ab)=\phi(a)\phi(b)$.
\end{enumerate}
A functor $\phi$ is an \emph{isomorphism} if it is bijective and $\phi^{-1}$ is also a functor.
\end{definition}

\begin{example}\label{examplepreordersarecategories}
Let $X$ be a set and $\leq$ be a preorder on $X$ (i.e., a reflexive and transitive relation), and $R=\left\{(y,x):x\leq y\right\}$ the associated subset of $X^2$. Let $\G$ be the directed graph associated with this pair $(X,R)$ as in Example \ref{examplegraphtriangle}. In this case, we have $\G^{(2)}=\left\{((z,y),(y,x)):(y,x),(z,y)\in R\right\}$, so defining composition by $(z,y)(y,x)=(z,x)$ we obtain a category structure on $\G$.

In fact, since $R$ is the set of arrows of $\G$ we identify $\G$ and $R$, and moreover, $R^{(0)}=\left\{(x,x):x\in X\right\}$ can be obviously identified with $X$ via $x\mapsto (x,x)$.

The categories $\mathcal{C}$ obtained in this manner (up to isomorphism) are precisely those such that, for every $x,y\in\mathcal{C}^{(0)}$, $\mathcal{C}_x^y$ has at most one arrow.
\end{example}

\begin{example}\label{examplemonoidsarecategories}
Let $M$ be a monoid (see Definition~\ref{definitioninversesemigroup}). Let $M^{(0)}=\left\{\ast\right\}$ be a singleton set, $M^{(1)}=M$ and $\so(m)=\ra(m)=\ast$ for every $m\in M$. In this case $M^{(2)}=M\times M$, and define composition as the monoid operation of $M$. This allows us to see $M$ as a category.

Conversely, if $\mathcal{C}$ is any category and $x\in\mathcal{C}^{(0)}$, then $\mathcal{C}_x^x$ is a monoid. If $\mathcal{C}^{(0)}=\left\{x\right\}$ is a singleton, then $\mathcal{C}$ is the category associated with the monoid $\mathcal{C}_x^x$ as above.
\end{example}

\begin{example}\label{monoidactioncategory}
Let $\theta$ be a global action (see Definition \ref{definitionactionspartialandglobal}) of a monoid $M$ on a set $X$, denoted by concatenation. We define the category $M\ltimes_\theta X$ whose underlying graph is similar to the one as in Example \ref{examplegraphdynamicalsystem}.

Namely, let $(M\ltimes_\theta X)^{(0)}=X$, $(M\ltimes_\theta X)^{(1)}=M\times X$, $\mathfrak{s}(m,x)=x$, $\mathfrak{r}(m,x)=mx$. In this case, $(M\ltimes_\theta X)^{(2)}=\left\{((m,nx),(n,x)):x\in X,m,n\in M\right\}$, so we define the composition as $(m,nx)(n,x)=(mn,x)$.

Note that if $X$ is a singleton $X=\left\{*\right\}$ and $M$ acts on $X$ trivially, we recover Example~\ref{examplemonoidsarecategories}.
\end{example}

We say that an arrow $a$ of a category $\mathcal{C}$ is \emph{invertible} if there exists another arrow $b$ such that $(a,b),(b,a)\in \mathcal{C}^{(2)}$, $ba=\mathfrak{s}(a)$ and $ab=\mathfrak{r}(a)$. In this case, the arrow $b$ is unique and it is denoted by $a^{-1}$, we call $a^{-1}$ the \emph{inverse} of $a$, and $a$ is called \emph{invertible}.

It is standard fact that $(a^{-1})^{-1}=a$ and $(ab)^{-1}=b^{-1}a^{-1}$ whenever $a$ and $b$ are invertible and $(a,b)\in\mathcal{C}^{(2)}$.

\begin{definition}[\cite{MR1724106}]
A \emph{groupoid} $\G$ is a small category in which all arrows are invertible.
\end{definition}

\begin{example}
Every group is a groupoid, in the same way as Example \ref{examplemonoidsarecategories}.
\end{example}

\begin{example}\nomenclature[G]{$G\ltimes_\theta X$}{Transformation groupoid}\index{Groupoid!Transformation groupoid}\label{exampletransformationgroupoid}
If $\theta$ is an action of a group $G$ on a set $X$ (as usual, denoted by concatenation), then the category $G\ltimes_\theta X$ (as in Example \ref{monoidactioncategory}) is a groupoid, which we call a \emph{transformation groupoid}. The inverse of an arrow $(g,x)\in G\ltimes_\theta X$ is $(g^{-1},gx)$.
\[\begin{tikzcd}[row sep=tiny]
x\arrow[r,"g"]\arrow[rr,bend right,"hg",swap]&g x\arrow[r,"h"]&(hg)x
\end{tikzcd}\]
When there is no risk of confusion, we will drop the index $\theta$ and denote this groupoid simply as $G\ltimes X$.
\end{example}

\begin{example}\label{exampleequivalencerelationsarecategories}
A preorder $R$ on a set $X$, regarded as a category as in Example \ref{examplepreordersarecategories}, is a groupoid if and only if $R$ is an equivalence relation. (See \ref{definitionprincipalgroupoid} below.)
\end{example}

\begin{example}\label{exampleorbitgroupoid}\index{Orbit!equivalence relation}\index{Orbit}\nomenclature[G]{$\mathcal{R}(\G)$}{Orbit equivalence relation of $\G$}\nomenclature[G]{$\mathcal{R}(\alpha)$}{Orbit groupoid of a group action $\alpha$}
If $\G$ is a groupoid, then the image of the map $(\ra,\so):\G\to\G[0]\times\G[0]$, $(\ra,\so)(a)=(\ra(a),\so(a))$, is an equivalence relation on $\G[0]$, called the \emph{orbit equivalence relation}, which we will denote by $\mathcal{R}(\G)$. An $\mathcal{R}(\G)$-equivalence class is called an \emph{orbit}.% In other words, two elements $x,y\in\G[0]$ belong to the same orbit if and only if there is $a\in\G$ with $\so(a)=x$ and $\ra(a)=y$.

\begin{itemize}
\item If $\theta$ is an action of a group $G$ on a space $X$, then $\mathcal{R}(G\ltimes_\theta X)$ is the usual orbit equivalence relation of the action $\theta$, so orbits of $G\ltimes_\theta X$ coincide with orbits of $\theta$. We will abuse notation and denote, in this case, the \emph{orbit groupoid} as $\mathcal{R}(\theta)=\mathcal{R}(G\ltimes_\theta X)$.
\item If $R$ is an equivalence relation on a set $X$, then the orbit of an element $x\in X$ is the $R$-equivalence class of $x$.
\end{itemize}
\end{example}

In fact, it is possible to determine the groupoids which are isomorphic to an equivalence relation, in the sense above, and those are called the \emph{principal groupoids}.

\begin{definition}\label{definitionisotropygroup}\index{Isotropy group}\nomenclature[G]{$\G_x^x$}{Isotropy group at $x$}
If $\G$ is a groupoid and $x\in\G[0]$, we call the set $\G_x^x$ the \emph{isotropy group} of $\G$ at $x$.
\end{definition}

Note that $\G_x^x$ is indeed a group with the operation inherited from $\G$.

\begin{definition}\index{Groupoid!Principal}\label{definitionprincipalgroupoid}
A groupoid $\G$ is \emph{principal} if $\G_x^x$ is a trivial group for all $x\in\G[0]$.
\end{definition}

\begin{example}\label{examplewhentransformationgroupoidisprincipal}
If $\theta$ is an action of a group $G$ on a set $X$, then the transformation groupoid $G\ltimes_\theta X$ is principal if, and only if, the action $\theta$ is free: Indeed, for all $x\in X$, let $G_x=\left\{g\in G:gx=x\right\}$ be the stabilizer of $x$, so the map
\[G_x\to (G\ltimes_\theta X)_x^x,\qquad g\mapsto (g,x)\]
is a group isomorphism. Thus $G\ltimes_\theta X$ is principal if and only if $G_x$ is trivial for all $x\in X$, i.e., $\theta$ is free.
\end{example}

\begin{proposition}\label{principalequalequivalencerelation}
A groupoid $\G$ is principal if and only if there exists an equivalence relation $R$ on a set $X$ which is isomorphic to $\G$, as a groupoid.
\end{proposition}
\begin{proof}
Every equivalence relation, regarded as a groupoid, is principal, and conversely if $\G$ is a principal groupoid then the map $(\ra,\so):\G\to\mathcal{R}(\G)$ is a groupoid isomorphism.\qedhere
\end{proof}

In the following sections and chapters, we will consider topological versions of these constructions, and generalize the transformation groupoids to \emph{groupoids of germs} of partial actions of inverse semigroups.

\subsection{Groupoids as algebraic structures}

A groupoid can be alternatively (algebraically) defined as a set $\G$ endowed with a partial binary map $\G[2]\to\G$, $(a,b)\mapsto ab$, where $\G[2]$ is a subset of the product $\G\times\G$, satisfying
\begin{enumerate}
\item[(G1)] If $(a,b)$ and $(b,c)\in\G[2]$, then $(ab,c)$ and $(a,bc)\in\G[2]$ and $(ab)c=a(bc)$;
\item[(G2)] For all $a\in\G$, there is $p\in\G$ such that $(p,a)\in\G[2]$ and if $(a,b)\in\G[2]$, then $(pa)b=b$;
\item[(G3)] For all $a\in\G$, there is $q\in\G$ such that $(a,q)\in\G[2]$ and if $(c,a)\in\G[2]$, then $c(aq)=c$.
\end{enumerate}
(note that the equations at the end of (G2) and (G3) make sense by (G1)).

Indeed, any groupoid, in the categorical sense, already comes with the structure described above. Conversely, if $\G$ is a groupoid as described above, the elements $p$ and $q$ in (G2) and (G3) are unique with those properties, and in fact they are equal so we denote $p=q=a^{-1}$. Define the source and range maps as $\so(a)=a^{-1}a$ and $\ra(a)=aa^{-1}$, and let $\G[0]=\ra(\G)=\so(\G)$, $\G[1]=\G$. This gives us all the remaining graph-theoretical structure.

Therefore, in order to describe a groupoid $\G$, it is enough to simply describe a partial binary operation on $\G$, which we identify with $\G$ as a set, satisfying (G1)-(G3) above, so the remaining structure is uniquely determined.

\begin{definition}\index{Morphism (of groupoids)}
A map $\phi:\G\to\H$ between two groupoids $\G$ and $\H$ is a \emph{morphism} if $(a,b)\in\G[2]$ implies $(\phi(a),\phi(b))\in\H[2]$, and $\phi(ab)=\phi(a)\phi(b)$.
\end{definition}

\begin{proposition}\label{propositiongroupoidmorphism}\index{Isomorphism (of groupoids)}
Let $\phi:\G\to\H$ be a groupoid morphism. Then
\begin{enumerate}
\item[(a)] $\phi(a^{-1})=\phi(a)^{-1}$ for all $a\in\G$;
\item[(b)] $\phi(\G[0])\subseteq\H[0]$;
\item[(c)] $\phi(\mathfrak{s}(a))=\mathfrak{s}(\phi(a))$ and $\phi(\mathfrak{r}(a))=\mathfrak{r}(\phi(a))$ for all $a\in\G$;
\item[(d)] If $\phi$ is invertible then $\phi^{-1}$ is also a morphism.
\end{enumerate}
(In particular groupoids morphisms are precisely their functors.)
\end{proposition}
\begin{proof}
\begin{enumerate}
\item[(a)] Since $\phi(a)=\phi(a)\phi(a^{-1})\phi(a)$, we can multiply both sides on the left and on the right by $\phi(a)^{-1}$ and obtain $\phi(a)^{-1}=\phi(a^{-1})$.
\item[(b)] and (c) follow from (a).
\item[(d)] Suppose $\phi$ is invertible and $(\phi(a),\phi(b))\in\H[2]$. From (c),
\[\phi(\mathfrak{s}(a))=\mathfrak{s}(\phi(a))=\mathfrak{r}(\phi(b))=\phi(\mathfrak{r}(b))\]
so $\mathfrak{s}(a)=\mathfrak{r}(b)$ and thus $\phi^{-1}(\phi(a)\phi(b))=\phi^{-1}(\phi(ab))=ab=\phi^{-1}(\phi(a))\phi^{-1}(\phi(b))$, therefore $\phi^{-1}$ is a morphism.\qedhere
\end{enumerate}
\end{proof}

\begin{example}\label{exampleunitgroupoid}\index{Groupoid!Unit groupoid}
Every set $X$ can be regarded as a groupoid by setting $X^{(2)}=\left\{(x,x):x\in X\right\}$ and the product $xx=x$ for all $x\in X$.

Alternatively, let $X^{(2)}$ be the identity equivalence relation on $X$, regarded as a groupoid as in \ref{exampleequivalencerelationsarecategories}, or let $\left\{1\right\}$ be the trivial group acting (trivially) on $X$ and consider the transformation groupoid $\left\{1\right\}\ltimes X$ as in \ref{exampletransformationgroupoid}. Any of the two bijections
\[\begin{array}{rcl}
X&\longrightarrow&X^{(2)}\\
x&\longmapsto&(x,x)
\end{array}\qquad\text{or}\qquad\begin{array}{rcl}
X&\longrightarrow&\left\{1\right\}\ltimes X\\
x&\longmapsto&(1,x)
\end{array}
\]
induces the same groupoid structure on $X$ as above.

In any case, the groupoids constructed in this manner are precisely the groupoids $\G$ satisfying $\G[0]=\G$, and these are called the \emph{unit groupoids}.
\end{example}

The next two examples show that the groupoid structure of a transformation groupoid is not enough to recover the initial group action.

\begin{example}
Let $G$ be a group acting on two sets $X$ and $Y$. We say that the $G$-spaces $X$ and $Y$ are \emph{conjugate} if there exist a bijection $F:X\to Y$ satisfying $F(gx)=g F(x)$ for all $g\in G$ and $x\in X$.

Clearly, if $X$ and $Y$ are conjugate, then the groupoids $G\ltimes X$ and $G\ltimes Y$ are isomorphic -- namely, $(g,x)\mapsto(g, F(x))$ is an isomorphism -- however the converse is not true.

Let $G=(\mathbb{Z}/2\mathbb{Z})^2=\langle a,b|a^2=b^2=[a,b]=1\rangle$ and $X=\left\{1,2,3,4\right\}$. We define two actions  $\alpha$ and $\beta$ of $G$ on $X$ on the generators in the following way:
\[\alpha(a)=(1\ 2),\quad\alpha(b)=(3\ 4),\quad\beta(a)=(1\ 2)(3\ 4)\quad\text{and}\quad\beta(b)=\operatorname{id}_X,\]
where $(i_1\cdots i_n)$ is denotes the cycle taking $i_j\mapsto i_{j+1}$. These actions are not conjugate, because $\alpha(a)$ is faithful but $\beta$ is not (even more strongly, there is no automorphism $\zeta$ of $G$ for which $\alpha$ is conjugate to $\beta\zeta$), but the respective transformation groupoids are isomorphic, and the isomorphism can be deduced from the usual depictions of these groupoids:
\[G\ltimes_\alpha X:\begin{tikzcd}
\bullet\arrow[loop left,"b"]\arrow[r,bend left=15,"a"]\arrow[r,bend left=50,"ab"]
	&\bullet\arrow[loop right,"b"]\arrow[l,bend left=15,"a"]\arrow[l,bend left=50,"ab"]
&		&\bullet\arrow[loop left,"a"]\arrow[r,bend left=15,"b"]\arrow[r,bend left=50,"ab"]
			&\bullet\arrow[loop right,"a"]\arrow[l,bend left=15,"b"]\arrow[l,bend left=50,"ab"]
\end{tikzcd}\]
\[G\ltimes_\beta X:\begin{tikzcd}
\bullet\arrow[loop left,"b"]\arrow[r,bend left=15,"a"]\arrow[r,bend left=50,"ab"]
	&\bullet\arrow[loop right,"b"]\arrow[l,bend left=15,"a"]\arrow[l,bend left=50,"ab"]
&		&\bullet\arrow[loop left,"b"]\arrow[r,bend left=15,"a"]\arrow[r,bend left=50,"ab"]
			&\bullet\arrow[loop right,"b"]\arrow[l,bend left=15,"a"]\arrow[l,bend left=50,"ab"]
\end{tikzcd}\]
\end{example}

\begin{example}
Following the previous example, suppose that two groups $G_1$ and $G_2$ act on spaces $X_1$ and $X_2$, respectively, and that the groupoids $G_1\ltimes X_1$ and $G_2\ltimes X_2$ are isomorphic. This clearly implies that $X_1$ and $X_2$ are in bijection, since they are the unit spaces of the respective groupoids, but it is not necessary that $G_1$ and $G_2$ are isomorphic.

Let $X=\left\{x,y\right\}$,
\[G_1=\mathbb{Z}/4\mathbb{Z}=\langle g|g^4=1\rangle\quad\text{and}\quad G_2=(\mathbb{Z}/2\mathbb{Z})^2=\langle a,b|a^2=b^2=[a,b]=1\rangle,\]
and consider actions of $\alpha_1$ of $G_1$ and $\alpha_2$ of $G_2$ on $X$ as
\[\alpha_1(g)=(x\ y),\qquad \alpha_2(a)=(x\ y),\qquad \alpha_2(b)=\operatorname{id}_X.\]
The map $\Phi:G_1\ltimes X\to G_2\ltimes X$ given by
\[\Phi(1,x)=(1,x),\quad\Phi(g,x)=(a,x),\quad\Phi(g^2,x)=(b,x),\quad\Phi(g^3,x)=(ab,x),\]
\[\Phi(1,y)=(1,y),\quad\Phi(g,y)=(a,y),\quad\Phi(g^2,y)=(b,y),\quad\Phi(g^3,y)=(ab,y)\]
is an isomorphism, even though $G_1$ and $G_2$ are not isomorphic.
\[G_1\ltimes_\beta X:\begin{tikzcd}
\bullet\arrow[loop left,"g^2"]\arrow[r,bend left=15,"g"]\arrow[r,bend left=50,"g^3"]
	&\bullet\arrow[loop right,"g^2"]\arrow[l,bend left=15,"g"]\arrow[l,bend left=50,"g^3"]
\end{tikzcd},\qquad
G_2\ltimes_\beta X:\begin{tikzcd}
\bullet\arrow[loop left,"b"]\arrow[r,bend left=15,"a"]\arrow[r,bend left=50,"ab"]
	&\bullet\arrow[loop right,"b"]\arrow[l,bend left=15,"a"]\arrow[l,bend left=50,"ab"]
\end{tikzcd}\]
\end{example}
\subsection{Constructions with groupoids}

\begin{definition}\nomenclature[G]{$\sqcup_i\G_i$}{Coproduct of groupoids}\label{definitiongroupoidcoproduct}
Let $\left\{\G_i:i\in I\right\}$ be a family of groupoids. We endow the disjoint union $\G=\sqcup_i \G_i$ with the set of composable pairs $\G[2]=\sqcup_i \G_i^{(2)}$ and the operations which restrict to the initial operation in each $\G_i$. We call this the \emph{coproduct} of the groupoids $\G_i$.
\end{definition}

From the definition of the coproduct we can readily see that $\sqcup \G_i$ satisfies the following universal property: There are canonical morphisms $\iota_i:\G_i\hookrightarrow\G$ (namely, the inclusion maps) such that for every groupoid $\H$ and every family of morphisms $\phi_i:\G_i\to\H$, there exists a unique morphism $\Phi:\G\to\H$ for which $\phi_i=\Phi\circ\iota_i$ for all $i$.

\begin{definition}
The \emph{product} of a family $\left\{\G_i\right\}$ of groupoids is the usual Cartesian product $\prod_i\G_i$ endowed with the partial operation $(a_i)_{i\in I}(b_i)_{i\in I}=(a_ib_i)_{i\in I}$, which is defined when every term $a_ib_i$ is defined.
\end{definition}

The product $\prod_i\G_i$ also satisfies the usual universal property: There are canonical morphisms $\pi_j:\prod_i\G_i\to\G_j$ (namely, the $j$-th projection map) such that for every groupoid $\H$ and family of morphisms $\phi_j:\H\to\G_j$, there exists a unique morphism $\Phi:\H\to\prod_i\G_i$ such that $\phi_j=\pi_j\circ \Phi$ for all $j$.

%\begin{definition}
%The \emph{orbit} of a unit $x\in\G[0]$ is the set $\ra(\G_x)=\left\{\ra(a):\so(a)=x\right\}$.
%\end{definition}

%Note that the distinct orbits of the units of $\G$ form a partition of $\G[0]$, which in turn defines a natural equivalence relation on $\G[0]$: $x$ and $y$ are equivalent if and only if $\G_x^y$ is nonempty. As a subset of $\G[0]\times\G[0]$, this equivalence relation is precisely the image of the map $(\ra,\so):\G\to\G[0]\times\G[0]$.

\begin{definition}[\cite{MR1712872,MR0244443}]
A groupoid is \emph{connected} or \emph{transitive} if it has only one orbit (Example \ref{exampleorbitgroupoid}), or equivalently if the map $(\ra,\so):\G\to\G[0]\times\G[0]$ is surjective.
\end{definition}

%\begin{example}
%If $G$ is a group acting on a set $X$, then the orbit of an element $x\in X$ inside the transformation groupoid $G\ltimes X$ is precisely the usual orbit $\left\{gx:g\in G\right\}$ given by the action of $G$. The groupoid $G\ltimes X$ is connected if and only if $G$ acts transitively on $X$.
%\end{example}

\begin{example}
The only connected equivalence relation on a set $X$ is $X^2$.
\end{example}

\begin{example}
If $\theta$ is an action of a group $G$ on a set $X$, then $G\ltimes_\theta X$ is connected if and only if $\theta$ is transitive.
\end{example}

\begin{proposition}\label{propositioneverygroupoidisacoproductofconnected}
Every groupoid $\G$ is a coproduct of connected groupoids.
\end{proposition}
\begin{proof}
Let us denote by $\operatorname{orb}(\G)$ the collection of orbits of elements of $\G[0]$. If $A\in\operatorname{orb}(\G)$ then the set $\G_A^A:=\so^{-1}(A)\cap\ra^{-1}(A)$ is a connected subgroupoid of $\G$, and it should be clear that $\G=\sqcup_{A\in\operatorname{orb}(\G)}\G_A^A$.\qedhere
\end{proof}

The next theorem shows that transformation groupoids - or groups and principal groupoids - can be seen as the ``building blocks'' for general groupoids.

\begin{theorem}\label{theoremclassificationofgroupoids}
\begin{enumerate}
\item[(a)] Every connected groupoid $\G$ is isomorphic to one of the form $G\times(X^2)$, where $G$ is a group, $X$ is a set and $X^2$ is the coarsest equivalence relation on $X$ as in Example \ref{exampleequivalencerelationsarecategories}.
\item[(b)] Every connected groupoid $\G$ is isomorphic to a transformation groupoid (of a transitive group action).
\item[(c)] Every groupoid $\G$ is isomorphic to a coproduct of transformation groupoids (of transitive group actions).
%\item[(d)] Every finite groupoid $G$ is isomorphic to a subgroupoid of a transformation groupoid of the form $S_N\ltimes[N]$.
\end{enumerate}
\end{theorem}
\begin{proof}
\begin{enumerate}
\item[(a)] Let $v\in\G[0]$ be an arbitrary element, and let $G=\G_v^v$, which is a group with the operation endowed from $\G$. For every $x\in\G[0]$, choose $t_x\in\G_v^x$.% Given $x,y\in\G[0]$, define the bijection $\phi_x^y:\G_x^y\to \G_v^v$, $\phi_x^y(a)=t_y^{-1}at_x$.

The map $G\times(\G[0])^2\to\G$ given by $(g,(y,x))\mapsto t_yg t_x^{-1}$ is a groupoid isomorphism, and its inverse is $a\mapsto(t_{\ra(a)}^{-1}at_{\so(a)},(\mathfrak{r}(a),\mathfrak{s}(a)))$.

\item[(b)] By item (a) we just need to consider a groupoid of the form $\G=G\times(X^2)$, where $G$ is a group and $X$ is a set. Endow $X$ with a group structure\ftnt{The statement that every set $X$ can be endowed with a group structure (in fact, abelian) is equivalent to the Axiom of Choice. If $X$ is finite this is clear enough, and if $X$ is infinite then it has the same cardinality as the set of its finite subsets, which is an abelian group under symmetric difference. For the converse, see \cite{MR0329895}.} and consider the action of $G\times X$ on $X$ as $(g,x)y=xy$.

Then the map $(G\times X)\ltimes X\to G\times(X^2)$, $((g,x),y)\mapsto(g,(xy,y))$ is a groupoid isomorphism, and its inverse is $(g,(x,y))\mapsto\left((g,xy^{-1}),y\right)$
\item[(c)] follows immediately from item (b) and the previous proposition.\qedhere
%\item[(c)] First note that $(S_N\ltimes[N])\dot{\cup}(S_M\ltimes[M])$ embeds naturally in $S_{N+M}\ltimes[N+M]$, by letting $S_N$ act only on the first elements of $[N+M]$ and $S_M$ act only on the remaining $M$ elements of $[N+M]$, with natural identifications. By the previous items, it suffices to consider a transformation groupoid $G\ltimes X$, where both $G$ and $X$ are finite. The action of $G$ on $X$ gives us a homomorphism $\phi:G\to S_X$. Let $\psi:G\to S_m$ be an injective homomorphism. For every $g\in G$, the disjoint union $\phi(g)\dot{\cup}\psi(g)$ is a bijection on $X\dot{\cup}[m]$. The map
%\[G\ltimes X\to S_{X\dot{\cup}[m]}\ltimes(X\dot{\cup}[m]),\qquad(g,x)\mapsto(\phi(g)\dot{\cup}\psi(g),x)\]
%is an injective homomorphism, and since $S_{X\dot{\cup}[m]}\ltimes(X\dot{\cup}[m])\simeq S_{[|X|+m]}\ltimes[|X|+m]$, we are done.\qedhere
\end{enumerate}
\end{proof}


%\end{document}